\documentclass[a4paper,10pt]{article}

\usepackage{textcomp}
\usepackage{enumitem}
\usepackage{amsmath}
\usepackage[a4paper, total={7in, 10in}]{geometry}
\usepackage{titlesec}
\usepackage{fancyhdr}
% \usepackage[sfdefault]{atkinson} %% Option 'sfdefault' if the base
%% font of the document is to be sans serif.
\usepackage[T1]{fontenc}
\usepackage{hyperref}
\usepackage{enumitem}
\usepackage{svg}
\usepackage[export]{adjustbox}
\usepackage{contour}
\usepackage{ulem}
\usepackage{listings}
\usepackage{xcolor}
\usepackage{minipage-marginpar}

\renewcommand{\ULdepth}{1.8pt}
\contourlength{0.8pt}

\lstdefinestyle{sparqlstyle}{
    basicstyle=\ttfamily\small,
    keywordstyle=\color{blue}\bfseries,
    commentstyle=\color{gray}\itshape,
    stringstyle=\color{orange},
    breaklines=true,
    frame=single,
    rulecolor=\color{black!30},
    backgroundcolor=\color{gray!8},
    captionpos=b,
    morekeywords={SELECT, WHERE, VALUES, ORDER, BY, PREFIX},
}

\lstdefinestyle{sqlstyle}{
    basicstyle=\ttfamily\small,
    keywordstyle=\color{blue}\bfseries,
    commentstyle=\color{gray}\itshape,
    stringstyle=\color{orange},
    breaklines=true,
    frame=single,
    rulecolor=\color{black!30},
    backgroundcolor=\color{gray!8},
    captionpos=b,
    morecomment=[l]{\#},
    morekeywords={SELECT, FROM, JOIN, ON, WHERE, ORDER, BY, ASC},
}

\newcommand{\myuline}[1]{%
  \uline{\phantom{#1}}%
  \llap{\contour{white}{#1}}%
}

\graphicspath{{images/}}

\urlstyle{same}
\hypersetup{
  colorlinks=true,
  linkcolor=blue,
  filecolor=magenta,
  urlcolor=cyan,
  pdftitle={Data Management Project},
}

\pagestyle{fancy}
\fancyhf{}
% \cfoot{\footnotesize Autorizzo il trattamento dei miei dati personali ai sensi del Dlgs 196 del 30 giugno 2003 e dell’art. 13 GDPR (Regolamento UE 2016/679) ai fini della ricerca e selezione del personale.}
\renewcommand{\headrulewidth}{0pt}

\title{Data Management Project \\
  \large DBMS Comparison}

\author{Leonardo Idone, 2088501}

\begin{document}

\maketitle

\section{Introduction}

In this document, we will qualitatively and quantitatively analyze the
differences between mapping a dataset to the relational model and to a
graph database model, specifically RDF. The dataset considered is
IMDb's Non-Commercial dataset, composed of information regarding
movies, actors, ratings, directors, writers and other attributes. As
for technology chosen, we used Apache Jena Fuseki for the Triplestore
and PostgreSQL for the Relational Database.

\section{Methodology}
\subsection{IMDb's Non-Commercial Datasets}

The dataset chosen has numerous problematic columns presenting a
classic ETL problem that gets resolved during data ingestion in
different ways:

\begin{itemize}{}
  \item Many rows have null values, designated by the character
  \texttt{\textbackslash N}, which needs to be handled cleanly.
  \item Columns like \texttt{primaryProfession} and
  \texttt{knownForTitles} have an array-like structure that does not
  map cleanly to either of our data models.
  \item Columns like \texttt{isAdult} and \texttt{isOriginalTitle}
  present boolean values in incompatible formats (0/1 instead of
  true/false).
  \item Some data is inconsistent; for instance, the
  \texttt{characters} column, where all strings are written in a
  single‑element JSON array
\end{itemize}

Another problem is the volume of data, about 9 GiB of tsv files gzipped
to 1.8 GiB. To effectively import them both smart buffering and
efficient secondary memory usage need to be employed, together with
maintaining an acceptable complexity during the ETL process, ideally
$O(n)$.

\subsection{Relational mapping}

We will now focus on the reasoning behind the chosen tables and
mapping to the respective columns.

For the file \texttt{title.basics} we have chosen to create the
following relations:

\begin{itemize}[label={}, noitemsep, before=\ttfamily]
  \item titleType(\myuline{titleTypeId}, titleTypeName)
  \item titleBasics(\myuline{tconst}, primaryTitle, originalTitle, isAdult, startYear, endYear, runtimeMinutes, *titleTypeId)
  \item genre(\myuline{genreId}, genreName)
  \item titleBasics\_genre(*\myuline{tconst}, *\myuline{genreId})
\end{itemize}

The column \texttt{tconst} is a natural primary key for each
title. Given the repetitive nature of \texttt{titleType}, the
\texttt{tconst - titleType} one-to-many relationship sees the latter
promoted to a new entity. The same happens for \texttt{tconst -
  genres}, with the difference that, being a many-to-many relationship
it will need a junction table other than a dedicated \texttt{genre}
relation. The rest of the columns act as attributes of the
\texttt{titleBasics} relation.

For the file \texttt{name.basics} we have chosen this subsequent
relations:

\begin{itemize}[label={}, noitemsep, before=\ttfamily]
  \item nameBasics(\myuline{nconst}, primaryName, birthYear, deathYear)
  \item knownForTitles(*\myuline{nconst}, *\myuline{tconst})
  \item profession(\myuline{professionId}, professionName)
  \item profession\_nameBasics(*\myuline{professionId}, *\myuline{nconst})
\end{itemize}

A string \texttt{nconst} is akin to \texttt{tconst} but for names,
making it a natural primary key. The columns
\texttt{primaryProfession} and \texttt{knownForTitles} are
many-to-many relationships between a \texttt{name} and multiple
\texttt{profession}(s) and \texttt{title}(s), so we follow the same logic as
before: we create a relation to avoid duplicates and insert junction
tables.

The data in \texttt{title.akas} has this relational schema:

\begin{itemize}[label={}, noitemsep, before=\ttfamily]
  \item titleAkas(\myuline{akasId}, *tconst, ordering, title, isOriginalTitle, *regionId, *languageId)
  \item region(\myuline{regionId}, regionName)
  \item language(\myuline{languageId}, languageName)
  \item akaType(\myuline{akaTypeId}, akaTypeName)
  \item akaType\_titleAkas(*\myuline{akaTypeId}, *\myuline{akasId})
  \item attribute(\myuline{attributeId}, attributeText)
  \item attribute\_titleAkas(*\myuline{attributeId}, *\myuline{akasId})
\end{itemize}

We denote \textit{akas} as a way of calling alternative titles for a
movie; for simplicity we decided to introduce an \texttt{akasId} to
distinguish them. This is a many-to-one relationship to
\texttt{titleBasics}. We also have the one-to-many relationships with
\texttt{region} and \texttt{language}. And finally, two many-to-many
relationships with \texttt{attributes} and with a new entity called
\texttt{akaType}.

As for \texttt{title.principals} we implemented this schema:

\begin{itemize}[label={}, noitemsep, before=\ttfamily]
  \item titlePrincipals(\myuline{principalsId}, *tconst, ordering, *nconst, *categoryId, *jobId, *characterId)
  \item character(\myuline{characterId}, characterName)
  \item category(\myuline{categoryId}, categoryName)
  \item job(\myuline{jobId}, jobName)
\end{itemize}

Title principals track what each person has done for a certain
title. What a person (\texttt{name}) has done can change a lot between
actor, director, camera operator and other roles. To track these roles
we have two one-to-many relationships in \texttt{character} and
\texttt{category}. If the job also can be correlated with a character
(for example in the case of an actor) we use the one-to-many
relationship \texttt{character}.

The table \texttt{title.crew} is very similar to
\texttt{title.principals} but focuses on the roles of writers and
directors:

\begin{itemize}[label={}, noitemsep, before=\ttfamily]
  \item writer(*\myuline{nconst}, *\myuline{tconst})
  \item director(*\myuline{nconst}, *\myuline{tconst})
\end{itemize}

Since these are effectively two many‑to‑many relationships between
\texttt{titles} and \texttt{names} all we have to do is create two
junction tables between them.

The data in \texttt{title.episode} is instead about distinguishing
which \texttt{titles} are singular movies and which are part of a
series:

\begin{itemize}[label={}, noitemsep, before=\ttfamily]
  \item titleEpisode(\myuline{episodeId}, *tconst, *parentTconst, seasonNumber, episodeNumber)
\end{itemize}

As such we just need to create two one-to-many foreign keys on the
relation: one for the \texttt{title} we are representing, and one for
the \texttt{title} representing the whole series. We also track the
season and episode number. Again we add an \texttt{episodeId} for
simplicity.

Finally we map the \texttt{title.ratings}:

\begin{itemize}[label={}, noitemsep, before=\ttfamily]
  \item ratings(*tconst, averageRating, numVotes)
\end{itemize}

This maps perfectly to a separate relation \texttt{ratings} with a
foreign key on the considered \texttt{title} and attributes on the
number of votes and average rating.

To import the data to PostgreSQL we have chosen to use Python with the
\texttt{psycopg} library. This allows direct piping of buffered rows
to the RDBMS via COPY protocol, allowing fast transactional operations
run while temporarily disabling integrity constraints until the
transaction is finished. As another expedient to speed up the
ingestion of data we have opted to create \texttt{unlogged} tables,
which are not crash resistant but fine for the purpose of our
project. We also used \texttt{rapidgzip} to overcome the single
threaded, small buffering of data done by the builtin python
\texttt{Gzip} library. The PostgreSQL library Psycopg uses a default
buffer size of 64 KiB for its operations, which did not act as a
bottleneck for the transactions. We implemented the strategy pattern
on the various ways rows should be ingested by each different file and
table, opening one connection for each separate table, one file at a
time, committing the whole converted file at once. This approach,
while fast enough, could be problematic for databases that refuse
multiple concurrent connections from one host. A better way would be
to generate the non-transactional data to be loaded sequentially in
the database.

\begin{figure}
  \begin{centering}
    \includesvg[width=1\textwidth,inkscapelatex=false]{imdb.sql.svg}
    \caption[IMDb Relational Schema]{IMDb's Relational Schema (note:
      some column names have been renamed for implementation
      reasons).}
    \label{fig:IMDb Relational Schema}
  \end{centering}
\end{figure}


\subsection{Triplestore Mapping}

We will explain the rationale behind the approach used to build the
Triplestore mapping. Since RDF is based on a way to allow
interconnection of elements primarily in the web, and seeing how IMDb
already implements it on its website, we decided to mimic its
structure where possible. In particular, IMDb makes heavy use of
\texttt{schema.org} models. On our part, in addition to mimicking
IMDb's choices, where possible we decided to expand it with
appropriate predicates and objects present in the most adaptable
schema. Specifically, we used the \texttt{RDF:type} of
\texttt{schema.org/Movie} to describe the \texttt{titles},
\texttt{schema.org/Person} for the \texttt{names},
\texttt{schema.org/AggregateRating} for the \texttt{ratings}. While
IMDb's implementation of the schema lacked a node identifier for each
triple, we included the subject by simply specifying the correct IMDb
page for each node. For example, a \texttt{title} with a
\texttt{tconst} will be mapped to the self-imposed IMDb schema (based
on their website mapping) of
\texttt{www.imdb.com/title/<tconst>}. Same thing for
\texttt{schema.org/AggregateRating}, which will map to
\texttt{www.imdb.com/title/<tconst>/reviews}. For every single other
column seen in the dataset, these will map to predicates present in
the aforementioned schemas used by IMDb itself. Some notable
exceptions are language and region, which are not tracked in
conformance with W3C's system of country codes, and as such are not
specified when considering alternative titles (even though
\texttt{schema.org} provides schemas for specifying detailed
alternative names, just not with IMDb's country codes). IMDb doesn't
bother tracking these country codes in their RDF schema, since they
use it for the semantic web and not archival purposes. Apart from that
lost information, everything was accounted for by \texttt{schema.org}
types. One thing of note is that any null value we came across simply
was not added to the graph.

To import the data we did not have any access to fast transactional
interfaces to Apache Jena Fuseki, so to ingest it we had to extend the
previous approach implemented in the relational part of this etl
process in two parts:

\begin{enumerate}
  \item Read all the files row by row, and convert them to a file
  format for fast parsing. In our case we chose N-Triples. One
  possibly better alternative is RDF Binary, also called RDF Thrift.
  \item Import them in a non-transactional way to Apache Jena Fuseki.
\end{enumerate}

The first part required us to build in Java the same strategy pattern
implemented by our SQL solution, switching from transactionally
copying to tables to non-transactionally dumping constructed
N-Triples. To do so we leveraged Apache Jena's \texttt{StreamRDF}
interface for fast writing together with Java's builtin
\texttt{BufferedReader} and \texttt{GZIP*Stream}, with
\texttt{FastCSV} for quick parsing and a buffer of 64 KiB where
configurable (this didn't speed up the process much compared to the
default 8 KiB buffer, the \texttt{ext4} filesystem in use has a block
size of 4 KiB). For the row generation we used the \texttt{Node} and
\texttt{NodeFactory} API for low-level construction of triples. All
operations are single threaded since the performance is acceptable
even without concurrent processing. It should be trivial to speed it
up even more by processing multiple files at the same time.

The second part, given the number of 225 million tuples, required
leveraging the large scale data loading tool \texttt{tdb2.xloader}.
Created by Apache, it works orders of magnitude better than the
standard tool \texttt{tdb2.tdbloader}. Their difference is the way
they build their SPO, POS, OSP indexes for subjects, objects and
predicates; \texttt{tdb2.xloader} does them sequentially instead of in
parallel (SPO\textrightarrow POS\textrightarrow OSP), trying to
maximize buffer usage, cache hits and using N-1 thread for parallel
processing of a index at a time. This is a better approach for large
amounts of data.

\begin{figure}
  \begin{centering}
    \includesvg[width=\textwidth,inkscapelatex=false]{graph.svg}
    \caption[Excerpt of RDF Graph]{Excerpt of RDF Graph for tt0109382.}
    \label{fig:Excerpt of RDF Graph.}
  \end{centering}
\end{figure}

\section{Analysis}

The implementation of the two technologies was significantly
different. Excluding the technological differences in loading the
data, adapting the original data to the schemas was much more
difficult in a relational setting than in a graph one. This is for one
simple reason: relational schemas are very granular and detailed,
always allowing us to find a way to map existing information to it. As
such, for every problematic column we found there was an approach that
allowed us to insert it efficiently in our schema.

Graph databases based on RDF instead are very standardized. While
technically schema-less, international conventions on how to store and
share data are present and encouraged in their use. These are very
comprehensive and rarely lack any desired fields, speeding up the
engineering and analysis phase immensely. Still, they require
integrating a private, personal schema for unimplemented requirements,
such as the custom country codes and language codes used by IMDb. This
is faster and simpler than building everything from scratch.

Another important difference is the lack of native support to
implement integrity checks at the schema level. Unlike SQL, RDF does
not provide a native way to do this, and W3C has since fixed this flaw
by introducing SHACL, a validation protocol for RDF graphs. This is
completely decoupled from the triples themselves and as such is
effectively separate from the data storage. This is expected, since
RDF and graph databases were first and foremost designed for rapidly
changing, extensible models and not archival needs. In other words, we
are able to set constraints, but these are not directly correlated to
a schema if not via external documentation. A possible, if still
decoupled, way of partially solving this can be found in the RDFS W3C
standard (constraints and schemas would still be two separate,
independent entities, but fixed in stone).

Instead, SQL allows for very precise integrity checks at the schema
level, declarable directly in the table definition. This enhances the
integrity of stored data since we can leverage both an external
validator and an internal checker.

Regarding the speed of queries, the simple model with which a graph
database is built allows for the construction of very thorough
triple-based indexes. These cover any association we may ever need
regarding subject-predicate-object queries, which are the vast
majority of query executed.

Conversely, SQL does not come with any indexes outside ones built for
every primary key in every table we have. In theory, we could map the
most common queries required by us and build what we need. But the
construction of multiple indexes can be costly and must be considered
carefully, making graph databases overall more versatile.

For example, comparing the same query for getting all the movies made
by one director, selected by director name, their title and their year
of publishing, both Apache Jena Fuseki and PostgreSQL spend similar
time to get a result, as long as the relational database has an index
on \texttt{nameBasics.primaryName}. As such, graph databases allow for
faster development and expansion of functionalities, since
implementing a new query for a service may not need any external
intervention from the engineer other than writing the SPARQL code.

\begin{figure}[h]
\begin{minipage}[t]{0.48\textwidth}
    \lstset{style=sparqlstyle}
    \begin{lstlisting}[caption={SPARQL Query}]
SELECT ?creativeWork ?directorName
       ?workName ?date
WHERE {
  VALUES ?directorName { "Nanni Moretti" }
  ?director schema:name ?directorName .
  ?creativeWork schema:director ?director ;
    schema:name ?workName ;
    schema:datePublished ?date .
} ORDER BY ?date
    \end{lstlisting}
\end{minipage}
\hfill
\begin{minipage}[t]{0.48\textwidth}
    \lstset{style=sqlstyle}
    \begin{lstlisting}[caption={SQL Query}]
-- Index on nb.primaryname suggested
SELECT 'https://www.imdb.com/title/'
       || tb.tconst,
       nb.primaryname,
       tb.primarytitle,
       tb.startyear
FROM namebasics nb
JOIN director d
  ON nb.nconst = d.nconst
JOIN titlebasics tb
  ON d.tconst = tb.tconst
WHERE nb.primaryname = 'Nanni Moretti'
ORDER BY tb.startYear ASC
    \end{lstlisting}
\end{minipage}
\end{figure}


\section{Conclusion}

Overall, we can say that graph databases, in this case specifically
RDF triplestores, provide a strong alternative to existing solutions
in the market. Unlike classical NoSQL solutions like MongoDB, graph
databases cover a much larger set of needs. While SQL is inherently
capable of addressing these needs, it requires significantly more
engineering. The only aspect that is objectively better in SQL is
integrity constraints, which are directly implemented in the schema
and effectively pair well with other client and server side validators
connected to the database. A more interesting research topic could be
the comparison of ORM tools for these two paradigms. Superficially
speaking, there do not seem to be particular hurdles that could limit
the usage of a graph database in an ORM environment, but more analysis
should be done.

\end{document}


%%% Local Variables:
%%% mode: LaTeX
%%% TeX-master: t
%%% TeX-command-extra-options: "--shell-escape"
%%% End:
